\chapter{The example and the two controls}\label{ch:example}

\begin{chapterplan}
The question of this chapter: what does the paper's one example look like, and what do its two controls say about
the source of the rise in entry?

After this chapter, you can do these things:
\begin{itemize}
\item Derive the six auction objects of the example in two ways, and derive the preparation threshold
  \(\tau(r)\) and its slope.
\item Show why, at the prior \(\tfrac12\), the hypotheses (H1)--(H3) of Proposition~2 reduce to one band for the
  top posterior, and state Corollary~1.
\item Explain the entry ceiling \(r_C\), and why the band of top posteriors is narrow.
\item State and prove the fixed-information control, Proposition~A.2, and say why its slot version needs a
  positive spread in each economy.
\item Read the two control panels of the paper's outcome table, and say what they show together.
\end{itemize}
\end{chapterplan}

% ====================================================================
\section{Why one example}\label{sec:example-why}

Chapters~\ref{ch:auction} to~\ref{ch:ccc} work with general classes and a general incumbent. Proposition~2 says: if
(H1), (H2) and (H3) hold, entry first rises and later falls across three incumbents. Two questions remain. Can the
three strict conditions hold at the same time, when (H1) needs gross profit to stay high against the strong
incumbent and (H3) needs it to fall low against the very strong one? And what do the objects look like in a
figure? The paper answers both with one economy, \emph{the example} (Section~4.2 of the paper). The numbers and
figures of the paper come from it.

In words, the example is a small economy with every part of the model. Each class has one value: a challenger of
class \(H\) has the value \(h\), and a challenger of class \(L\) has the value \(\ell\). So the class
and the value coincide, and preparation teaches the challenger nothing beyond its class. The incumbent's value is
uniform on an interval \([0,r]\). One number, the strength \(r\), then describes the incumbent, and a larger \(r\)
is a stronger incumbent. With these choices every integral of Chapter~\ref{ch:auction} has a short closed form.

\begin{analogy}[a prototype]
An engineer has a list of requirements that pull against each other, and builds one prototype that meets all of
them. The prototype shows that the list is not contradictory and lets everyone see the parts. The example plays
this role for Proposition~2, and small changes of its parameters keep all three conditions
(Section~\ref{sec:example-open}).

Where the analogy breaks: a prototype is a physical test, and its success could be luck. Every statement about the
example is an exact consequence of the propositions.
\end{analogy}

The first part of the chapter computes the example. The second part uses it for two controls. Each control removes
one channel and asks whether the rise in entry survives.

% ====================================================================
\section{The primitives}\label{sec:example-primitives}

This section lists the example's primitives in symbols. Chapter~\ref{app:numbers} gives their values.

\begin{itemize}
\item \emph{Classes.} The class distribution \(G_H\) puts all its mass on \(h\), and \(G_L\) puts all its mass on
  \(\ell\). In the notation of Section~\ref{sec:toolkit-cdf}, \(G_H(y)=\1\{y\ge h\}\) and \(G_L(y)=\1\{y\ge\ell\}\).
\item \emph{Reserve.} The seller's reserve is \(p>0\).
\item \emph{Incumbent.} The incumbent's value \(R\) is uniform on \([0,r]\). Its distribution function is
  \(F_r(y)=y/r\) for \(0\le y\le r\) and \(F_r(y)=1\) for \(y\ge r\) (Section~\ref{sec:toolkit-cdf}).
\item \emph{Order of the values.} The values, the reserve and every strength on the path satisfy
  \begin{equation}\label{eq:example-order}
  0<p<\ell<r<h.
  \end{equation}
\item \emph{Other primitives.} The prior is \(\pi=\Pr(j=H)\), the zero-cost probability is \(\rho\), the high cost
  is \(c_H\), and the trading cost is \(k\). Money is measured in units of the low value. The weak, the strong and
  the very strong incumbent have the strengths \(r_0<r_1<r_2\). The paper fixes the prior at \(\tfrac12\).
\end{itemize}

Each inequality in \eqref{eq:example-order} has a job.
\begin{itemize}
\item \(p<\ell\): the reserve lies below both values, so every value meets the reserve. With \(\ell<h\), this gives
  the class ranking \eqref{eq:model-classes}: \(G_H\le G_L\) everywhere and \(G_H(p^-)=G_L(p^-)=0\).
\item \(\ell<r\): the incumbent's value lies above \(\ell\) with the positive probability \(1-\ell/r\). So the
  incumbent sometimes beats a class-\(L\) challenger, and Section~\ref{sec:example-auction} shows that this makes
  the spread positive.
\item \(r<h\): the incumbent's value never reaches \(h\). So a prepared class-\(H\) challenger always wins, and
  \(W_H(F_r)=F_r(h)=1\). Recall from \eqref{eq:model-W} that \(W_j(F)\) is the probability that a prepared class-\(j\)
  challenger beats the incumbent.
\end{itemize}

\emph{Why a larger \(r\) is a stronger incumbent.} For \(r'>r\), \(F_{r'}(y)\le F_r(y)\) at every \(y\): on
\([0,r]\), \(y/r'\le y/r\); on \((r,r']\), \(y/r'\le1\); above \(r'\), both equal one. So an increase in \(r\) is a
first-order stochastic strengthening, and the uniform laws form a path (Section~\ref{sec:toolkit-fosd}). Every
comparison along the path solves the equilibrium again at each strength.

\emph{The slot.} Recall the order slot. With probability \(\lambda\) the informed investor fills the slot. Otherwise
a liquidity trader fills it, and it sells, abstains or buys with the probabilities \(\beta\), \(1-2\beta\) and
\(\beta\). A liquidity order in one direction has probability \(b=(1-\lambda)\beta\), and a liquidity abstention
has probability \(b_0=(1-\lambda)(1-2\beta)\). The mass ratio is \(u=\lambda/b\), and the largest likelihood ratio of
one order is \(1+u\) (Chapter~\ref{ch:beliefs}). In the example's slot, \(1+u=\eu\), Euler's number, and \(\beta\)
lies inside the model's domain \(\beta\le1/3\) (Section~4.2 of the paper). Corollary~1 below lets the slot vary over
\(0<\lambda<1\) and \(0<\beta<1/2\). Keep the two domains apart.

% ====================================================================
\section{The auction in the example}\label{sec:example-auction}

The auction objects are \(t_0\), \(t_H\), \(t_L\), \(g_H\), \(g_L\) and \(\DT=t_H-t_L\). Recall what they mean
(Chapter~\ref{ch:auction}): \(t_0\) is the expected target proceeds when no challenger prepares; \(t_j\) is the
expected target proceeds after a class-\(j\) challenger prepares; \(g_j\) is that challenger's expected gross
acquisition profit, before its cost. This section derives all six in two ways.

\subsection{Route 1: the case table}

Each bidder bids its value (Section~\ref{sec:auction-bids}). So the winner pays the larger of the reserve and the
other bidder's value, when that bidder bids. Table~\ref{tab:example-cases} applies this rule to the example. A tie
\(R=\ell\) has probability zero, because \(F_r\) is continuous, so the table ignores it.

\begin{table}[htbp]
\centering\small
\caption{Who wins and who pays in the example, at each incumbent value \(R\in[0,r]\).}
\label{tab:example-cases}
\begin{tabular}{@{}lllll@{}}
\toprule
Incumbent value & Class \(H\): winner & price & Class \(L\): winner & price\\
\midrule
\(R<p\) & challenger & \(p\) & challenger & \(p\)\\
\(p\le R<\ell\) & challenger & \(R\) & challenger & \(R\)\\
\(\ell<R\le r\) & challenger & \(R\) & incumbent & \(\ell\)\\
\bottomrule
\end{tabular}
\end{table}

When \(R<p\), the incumbent does not bid, and either challenger buys at \(p\). When \(p\le R<\ell\), either
challenger wins and pays \(R\). When \(\ell<R\le r\), class \(H\) still wins, since \(R\le r<h\), while class \(L\)
loses and the incumbent pays its bid \(\ell\). The two classes differ only in the last row. A winner keeps its
value minus the price, and a loser keeps nothing.

Every expectation over \(R\) has the same form. The density of \(R\) is \(1/r\) on \([0,r]\), so
\(\E[\varphi(R)]=\frac1r\int_0^r\varphi(y)\,dy\) for any payment \(\varphi\). The integral adds up the payment at
each incumbent value, and the factor \(1/r\) turns the sum into an average.

Without a challenger, the incumbent buys at \(p\) when \(R\ge p\), so \eqref{eq:model-t0} gives
\(t_0=p\,[1-F_r(p)]=p(1-p/r)\), since \(F_r(p)=p/r\) for \(p<r\).

\begin{derivation}[the proceeds with a class-\(H\) challenger]
The target receives \(p\) when \(R<p\) and \(R\) when \(R\ge p\) (first column of prices).
\begin{align*}
t_H&=\frac1r\Bigl[\int_0^p p\,dy+\int_p^r y\,dy\Bigr] \why{the class-\(H\) prices}\\
&=\frac1r\Bigl[p\cdot p+\frac{r^2-p^2}{2}\Bigr] \why{constant times length; \(\int y\,dy=y^2/2\)}\\
&=\frac1r\cdot\frac{2p^2+r^2-p^2}{2} \why{common denominator}\\
&=\frac{r^2+p^2}{2r}=\frac r2+\frac{p^2}{2r} \why{split the fraction}
\end{align*}
\end{derivation}

In words, \(t_H=\E[\max\{p,R\}]\): class \(H\) always wins, and it pays the incumbent's value or the reserve,
whichever is larger.

\begin{derivation}[the proceeds with a class-\(L\) challenger]
The target receives \(p\), then \(R\), then \(\ell\) on the three rows.
\begin{align*}
t_L&=\frac1r\Bigl[\int_0^p p\,dy+\int_p^\ell y\,dy+\int_\ell^r \ell\,dy\Bigr] \why{the class-\(L\) prices}\\
&=\frac1r\Bigl[p^2+\frac{\ell^2-p^2}{2}+\ell(r-\ell)\Bigr] \why{the three integrals}\\
&=\frac1r\Bigl[\frac{p^2}{2}+\frac{\ell^2}{2}-\ell^2+\ell r\Bigr] \why{expand}\\
&=\frac1r\Bigl[\ell r-\frac{\ell^2-p^2}{2}\Bigr] \why{collect \(\ell^2\) and \(p^2\)}\\
&=\ell-\frac{\ell^2-p^2}{2r} \why{divide each term by \(r\)}
\end{align*}
\end{derivation}

Class \(H\) wins at every \(R\) and keeps \(h-\max\{p,R\}\). So
\begin{align*}
g_H&=\frac1r\int_0^r h\,dy-\frac1r\int_0^r\max\{p,y\}\,dy \why{the class-\(H\) rows; linearity}\\
&=h-t_H=h-\frac r2-\frac{p^2}{2r} \why{the second term is \(t_H\)}
\end{align*}

\begin{derivation}[the gross profit of class \(L\)]
Class \(L\) keeps \(\ell-p\) when \(R<p\), keeps \(\ell-R\) when \(p\le R<\ell\), and keeps nothing when it loses.
\begin{align*}
g_L&=\frac1r\Bigl[\int_0^p(\ell-p)\,dy+\int_p^\ell(\ell-y)\,dy+\int_\ell^r0\,dy\Bigr] \why{the class-\(L\) rows}\\
&=\frac1r\Bigl[p(\ell-p)+\frac{(\ell-p)^2}{2}\Bigr] \why{see the next line}\\
&=\frac{\ell-p}{r}\Bigl[p+\frac{\ell-p}{2}\Bigr] \why{factor out \(\ell-p\)}\\
&=\frac{\ell-p}{r}\cdot\frac{\ell+p}{2}=\frac{\ell^2-p^2}{2r} \why{\((\ell-p)(\ell+p)=\ell^2-p^2\)}
\end{align*}
The middle integral is
\(\int_p^\ell(\ell-y)\,dy=\bigl[\ell y-\tfrac{y^2}{2}\bigr]_p^\ell=\tfrac{\ell^2}{2}-\ell p+\tfrac{p^2}{2}
=\tfrac{(\ell-p)^2}{2}\).
\end{derivation}

Subtract to get the spread. The \(p^2\) terms cancel:
\[
\DT=t_H-t_L=\frac r2-\ell+\frac{\ell^2}{2r}=\frac{r^2-2\ell r+\ell^2}{2r}=\frac{(r-\ell)^2}{2r}.
\]

\subsection{Route 2: the general formulas}

Chapter~\ref{ch:auction} gives every auction object as an integral over levels \(y\). For the point masses,
\(1-G_H(y)=\1\{y<h\}\), \(1-G_L(y)=\1\{y<\ell\}\), and the class difference is
\(D(y)=G_L(y)-G_H(y)=\1\{\ell\le y<h\}\), with \(D(p^-)=0\) (Section~\ref{sec:auction-slices}).

\begin{derivation}[the gross profits from \eqref{eq:auction-gross}]
The integral adds up, over levels \(y\) above the reserve, the probability that the challenger's value passes
\(y\) times the probability that the incumbent lies below \(y\).
\begin{align*}
g_H&=\int_p^\infty\1\{y<h\}\,F_r(y)\,dy=\int_p^hF_r(y)\,dy \why{\eqref{eq:auction-gross}}\\
&=\int_p^r\frac yr\,dy+\int_r^h1\,dy \why{split at \(r\); \(p<r<h\)}\\
&=\frac{r^2-p^2}{2r}+(h-r) \why{the two integrals}\\
&=h-\frac r2-\frac{p^2}{2r} \why{\(\tfrac{r^2}{2r}-r=-\tfrac r2\)}\\[4pt]
g_L&=\int_p^\infty\1\{y<\ell\}\,F_r(y)\,dy=\int_p^\ell\frac yr\,dy \why{\eqref{eq:auction-gross}; \(\ell<r\)}\\
&=\frac{\ell^2-p^2}{2r} \why{\(\int y\,dy=y^2/2\)}
\end{align*}
\end{derivation}

\begin{derivation}[the spread from \eqref{eq:auction-spread}]
The integral adds up, over levels \(y\), the class difference \(D(y)\) times the probability that the incumbent
lies above \(y\). That is the chance that the slice at \(y\) reaches the target.
\begin{align*}
\DT&=pF_r(p)\,D(p^-)+\int_p^\infty\bigl[1-F_r(y)\bigr]D(y)\,dy \why{\eqref{eq:auction-spread}}\\
&=0+\int_\ell^h\bigl[1-F_r(y)\bigr]\,dy \why{\(D(p^-)=0\); \(D=\1\{\ell\le y<h\}\)}\\
&=\int_\ell^r\Bigl(1-\frac yr\Bigr)dy+\int_r^h0\,dy \why{\(F_r=1\) above \(r\)}\\
&=(r-\ell)-\frac{r^2-\ell^2}{2r} \why{the integral}\\
&=(r-\ell)\Bigl[1-\frac{r+\ell}{2r}\Bigr] \why{\(r^2-\ell^2=(r-\ell)(r+\ell)\)}\\
&=(r-\ell)\cdot\frac{r-\ell}{2r}=\frac{(r-\ell)^2}{2r} \why{\(1-\tfrac{r+\ell}{2r}=\tfrac{r-\ell}{2r}\)}
\end{align*}
\end{derivation}

The proceeds follow from the split \eqref{eq:auction-split}: for a value that meets the reserve,
\(t(\theta)+g(\theta)=\theta\). Each class has one value, so
\begin{align*}
t_H&=h-g_H=\frac r2+\frac{p^2}{2r} \why{split at \(\theta=h\)}\\
t_L&=\ell-g_L=\ell-\frac{\ell^2-p^2}{2r} \why{split at \(\theta=\ell\)}
\end{align*}
Both routes agree.

\subsection{A third look at the spread}

At one incumbent value, the class difference in proceeds is \(\Gamma(R)=(R-\ell)^+\) for \(R<h\)
(Section~\ref{sec:auction-slices}): the last row of the table. The spread is its average:
\begin{align*}
\DT&=\E\bigl[(R-\ell)^+\bigr] \why{\(\DT=\E_F[\Gamma(R)]\)}\\
&=\frac1r\int_\ell^r(y-\ell)\,dy \why{\((y-\ell)^+=0\) for \(y\le\ell\)}\\
&=\frac1r\Bigl[\frac{(y-\ell)^2}{2}\Bigr]_\ell^r=\frac{(r-\ell)^2}{2r} \why{the integral}
\end{align*}
This form splits into two factors: \(\DT=\Pr(R>\ell)\cdot\E[R-\ell\mid R>\ell]\). Given \(R>\ell\), the value \(R\) is
uniform on \((\ell,r]\), so the conditional mean of \(R-\ell\) is \((r-\ell)/2\). Hence
\[
\DT=\frac{r-\ell}{r}\cdot\frac{r-\ell}{2}.
\]

The six objects of the example, equation (8) of the paper (with the integrals in its equation (A.1)), are
\begin{equation}\label{eq:example-payoffs}
\begin{aligned}
t_0&=p\Bigl(1-\frac pr\Bigr),& t_H&=\frac r2+\frac{p^2}{2r},& t_L&=\ell-\frac{\ell^2-p^2}{2r},\\
g_H&=h-\frac r2-\frac{p^2}{2r},& g_L&=\frac{\ell^2-p^2}{2r},& \DT(r)&=\frac{(r-\ell)^2}{2r}.
\end{aligned}
\end{equation}
As a check, \(t_H+g_H=h\) and \(t_L+g_L=\ell\), as the split requires.

\begin{intuition}[what the six formulas say]
\begin{itemize}
\item Class \(H\) always wins, so its profit is its value minus its expected payment, \(g_H=h-t_H\).
\item Class \(L\) earns only when the incumbent lies below \(\ell\), and its margin \(\ell-\max\{p,R\}\) is at most
  \(\ell-p\). So \(g_L\) is small, and it shrinks like \(1/r\) as the incumbent strengthens.
\item The spread needs two things at once: the incumbent must beat class \(L\), and it must then pay a price
  \(\ell\) that falls short of what class \(H\) would pay. Each factor is small when \(r\) is near \(\ell\).
\end{itemize}
\end{intuition}

% ====================================================================
\section{How the objects move with strength}\label{sec:example-shape}

\subsection{The three derivatives}

Write \(\DT\) as a sum before you differentiate:
\(\DT(r)=\tfrac{r^2-2\ell r+\ell^2}{2r}=\tfrac r2-\ell+\tfrac{\ell^2}{2r}\). The derivative of \(1/r\) is
\(-1/r^2\).

\begin{derivation}[the derivatives of the spread and the two gross profits]
\begin{align*}
\DT'(r)&=\frac12-\frac{\ell^2}{2r^2} \why{differentiate \(\tfrac r2-\ell+\tfrac{\ell^2}{2r}\)}\\
&=\frac{r^2-\ell^2}{2r^2}>0 \why{\(r>\ell>0\)}\\[4pt]
g_H'(r)&=-\frac12+\frac{p^2}{2r^2} \why{differentiate \(h-\tfrac r2-\tfrac{p^2}{2r}\)}\\
&=-\frac{r^2-p^2}{2r^2}<0 \why{\(r>p>0\)}\\[4pt]
g_L'(r)&=-\frac{\ell^2-p^2}{2r^2}<0 \why{differentiate \(\tfrac{\ell^2-p^2}{2r}\); \(\ell>p\)}
\end{align*}
\end{derivation}

These are the paper's derivatives, equation (9) of Section~4.2. Their signs follow from \eqref{eq:example-order}:
\(r>\ell\), \(r>p\) and \(\ell>p\). Proposition~1 promised only weak directions for a general strengthening. In the
example all three changes are strict at every strength on the path.

Gross profit at a belief \(\mu\) is \(B_r(\mu)=(1-\mu)\,g_L(r)+\mu\,g_H(r)\). So
\begin{align*}
\frac{\partial B_r(\mu)}{\partial r}&=(1-\mu)\,g_L'(r)+\mu\,g_H'(r) \why{differentiate term by term}\\
&<0 \why{both derivatives \(<0\); weights \(\ge0\), sum one}
\end{align*}
for every \(\mu\in[0,1]\) (Appendix~A.1 of the paper). Gross profit falls strictly at every belief, while the
spread rises strictly. The two forces of Chapter~\ref{ch:ccc} move in opposite directions at every strength above
\(\ell\).

\emph{The payment of class \(H\).} Since \(g_H=h-t_H\), the fall of \(g_H\) is the rise of class \(H\)'s expected
payment \(t_H=r/2+p^2/(2r)\):
\[
t_H'(r)=\frac12-\frac{p^2}{2r^2},\qquad 0<t_H'(r)<\frac12 .
\]
Class \(H\) pays \(\max\{p,R\}\). The mean of a uniform value on \([0,r]\) is \(r/2\), so it rises by one half per
unit of strength. The reserve puts a floor under the payment, and this floor slows the rise a little. So class
\(H\)'s payment rises by less than one half per unit of strength. Section~\ref{sec:example-ceiling} uses this fact.

\subsection{The class-gap identity in the example}

The profit gap is
\begin{align*}
\gamma(r)=g_H-g_L&=h-\frac r2-\frac{p^2}{2r}-\frac{\ell^2-p^2}{2r} \why{\eqref{eq:example-payoffs}}\\
&=h-\frac r2-\frac{\ell^2}{2r} \why{the \(p^2\) terms cancel}\\
&=\frac{2rh-r^2-\ell^2}{2r} \why{common denominator \(2r\)}
\end{align*}
It is positive: Chapter~\ref{ch:auction} (Check yourself~\ref{q:auction-pointgross}) gives
\(g_H-g_L=\int_\ell^hF_r(y)\,dy\), and \(F_r(y)>0\) for \(y>0\).

The class-gap identity \eqref{eq:auction-gap} says that \(\DT+g_H-g_L\) equals
\(\E_H[\max\{\theta,p\}]-\E_L[\max\{\theta,p\}]\). In the example the right side is \(h-\ell\), since both values
exceed \(p\). Check the left side:
\begin{align*}
\DT+g_H-g_L&=\frac{(r-\ell)^2}{2r}+\frac{2rh-r^2-\ell^2}{2r} \why{the two formulas}\\
&=\frac{r^2-2r\ell+\ell^2+2rh-r^2-\ell^2}{2r} \why{expand \((r-\ell)^2\)}\\
&=\frac{2rh-2r\ell}{2r}=h-\ell \why{\(r^2\) and \(\ell^2\) cancel}
\end{align*}
The right side of the identity is \(h-\ell\) at every \(r\) (Section~4.2 of the paper). So every unit by which the
spread rises is a unit by which the profit gap falls. The derivatives agree:
\(g_H'-g_L'=-\tfrac12+\tfrac{\ell^2}{2r^2}=-\DT'\).

\subsection{The shape of the spread}

The formula \(\DT(r)=(r-\ell)^2/(2r)\) gives the shape.
\begin{itemize}
\item At \(r=\ell\), \(\DT=0\). Its slope there is also zero: \(\DT'(\ell)=\tfrac12-\tfrac{\ell^2}{2\ell^2}=0\).
\item Above \(\ell\), \(\DT'>0\), so the spread rises with strength.
\end{itemize}
The reason is in the two factors of Section~\ref{sec:example-auction}. When \(r\) is near \(\ell\), the incumbent
almost never outbids a class-\(L\) challenger: \(\Pr(R>\ell)=(r-\ell)/r\) is small. And when it does, its value
lies only a little above \(\ell\), so class \(H\) pays only a little more than the incumbent pays for class \(L\).
The product of two small factors is very small, and the spread starts flat (Section~4.2 of the paper).

\begin{checkyourself}
\label{q:example-reserve}
The spread \(\DT=(r-\ell)^2/(2r)\) contains no \(p\).
(a)~Use Table~\ref{tab:example-cases} to explain why the reserve does not affect the spread.
(b)~What happens to \(g_L\) as \(p\) rises toward \(\ell\), with \(p<\ell\) kept? Explain in words.
\end{checkyourself}

\subsection{The preparation threshold}

Recall the preparation threshold \eqref{eq:beliefs-threshold}: when \(g_H>g_L\), it is the belief at which gross
profit equals the high cost, \(\tau=(c_H-g_L)/(g_H-g_L)\). A high-cost challenger prepares exactly when its belief
reaches \(\tau\), under the default tie rule.

\begin{derivation}[the threshold of the example]
Multiply the numerator and the denominator by \(2r>0\).
\begin{align*}
\tau(r)&=\frac{c_H-g_L}{g_H-g_L} \why{\eqref{eq:beliefs-threshold}}\\
&=\frac{c_H-\frac{\ell^2-p^2}{2r}}{\frac{2rh-r^2-\ell^2}{2r}} \why{\(g_L\) and \(\gamma\) above}\\
&=\frac{2rc_H-(\ell^2-p^2)}{2rh-r^2-\ell^2} \why{multiply above and below by \(2r\)}
\end{align*}
\end{derivation}

This is the paper's threshold (Section~4.2):
\begin{equation}\label{eq:example-tau}
\tau(r)=\frac{2rc_H-(\ell^2-p^2)}{2rh-r^2-\ell^2}.
\end{equation}
The denominator equals \(2r\gamma(r)>0\).

\emph{The threshold rises with strength.} Lemma~A.5 (Section~\ref{sec:trading-path}) says that the threshold never
falls along a path where \(\gamma>0\) and \(\tau\ge0\). Its increment has the numerator
\(g_L(F_r)-g_L(F_{r'})+\tau(r)[\DT(r')-\DT(r)]\). Both conditions hold on the example's path:
\begin{itemize}
\item \(\gamma(r)>0\), shown above;
\item \(\tau\ge0\) means \(g_L\le c_H\); Section~\ref{sec:trading-path} shows that without it the threshold can
  fall. At the example's values, \(\tau(\ell)>\pi>0\) (Appendix~A.6 of the paper), so
  \(c_H>g_L(\ell)\). Since \(g_L\) falls with \(r\), \(c_H>g_L(r)\) at every \(r\ge\ell\).
\end{itemize}
In the example \(g_L\) falls strictly, so the first term of the numerator is positive and the second is
nonnegative. So \(\tau\) rises strictly on the path (Appendix~A.6 of the paper: \(\tau\) rises on \(\ell\le r\le h\)).

At \(r=\ell\) the denominator is \(2\ell(h-\ell)\), so \(\tau(\ell)=[2\ell c_H-(\ell^2-p^2)]/[2\ell(h-\ell)]\).
Section~\ref{sec:example-control} uses this value.

% ====================================================================
\section{Corollary 1: the band of slots}\label{sec:example-corollary}

Recall the three conditions of Proposition~2 (Chapter~\ref{ch:ccc}), with \(B_i=B_{r_i}\) and
\(\Delta_i=\DT(r_i)\):
\[
\text{(H1)}\ \ B_0(\pi)<c_H<B_1(M(\pi)),\qquad
\text{(H2)}\ \ \Delta_0<k<\rho\,\kappa(\pi)\,\Delta_1,\qquad
\text{(H3)}\ \ B_2(M(\pi))<c_H.
\]
Here \(M(\pi)\) and \(m(\pi)\) are the top and the bottom posterior \eqref{eq:beliefs-bounds}, and
\(\kappa(\pi)=\min\{m(\pi),1-M(\pi)\}\) is the factor in the trading bound \eqref{eq:trading-bound}. In the example,
the auction is fixed by \(h\), \(\ell\), \(p\) and \(r\), so only the slot and the prior remain free inside these
conditions.

For the statement, write three numbers built from the example's auction, costs and strengths:
\[
u_{\rm lo}=\frac{2\tau(r_1)-1}{1-\tau(r_1)},\qquad
u_{\rm hi}=\frac{2\tau(r_2)-1}{1-\tau(r_2)},\qquad
u_k=\frac{\rho\,\DT(r_1)}{k}-2.
\]
The paper prints their values. Here they are the band of Chapter~\ref{app:numbers}.

\needspace{12\baselineskip}
\begin{paperresult}{Corollary 1 (Recruitment and ownership in the example, full statement)}
Let the classes be point masses at \(h\) and \(\ell\), the incumbent be uniform on \([0,r]\), and let
\((p,\rho,c_H,k)\) and \((r_0,r_1,r_2)\) take the example's values. Take a slot with \(0<\lambda<1\),
\(0<\beta<1/2\), \(u=\lambda/[(1-\lambda)\beta]\).
\begin{enumerate}[label=(\alph*),leftmargin=2em]
\item At prior \(\pi=\tfrac12\), the high-cost window, trading-cost window and exclusion condition hold together
  exactly on \(u_{\rm lo}<u<u_{\rm hi}\), equivalently \(\tau(r_1)<M(\pi)<\tau(r_2)\), with the endpoints of the band
  of Chapter~\ref{app:numbers}. This is an interval of presence probabilities at each fixed liquidity share. Under
  all four tie rules, entry is \(\rho\), \(\rho+(1-\rho)(b+\lambda/2)\), \(\rho\) at the three strengths, and
  high-class ownership is \(\rho/2\), \([\rho+(1-\rho)(b+\lambda)]/2\), \(\rho/2\).
\item For \(u<u_{\rm lo}\), entry at \(r_1\) is \(\rho\). At \(u=u_{\rm lo}\), orders are uniquely zero at \(r_0\)
  and uniquely full at \(r_1,r_2\). When an indifferent high-cost challenger prepares, every conclusion of
  Proposition~2 holds. When it may prepare with any probability, entry after a purchase at \(r_1\) ranges over
  \([\rho,1]\), and total entry there ranges over \([\rho,\rho+(1-\rho)(b+\lambda/2)]\). For
  \(u_{\rm hi}\le u<u_k\), a purchase at \(r_2\) induces certain preparation when an indifferent high-cost
  challenger prepares. For \(u\ge u_k\), the strict upper inequality of the trading-cost window at \(r_1\) fails.
\item If \(1+u=\eu\), all three hypotheses hold exactly when
  \[
  \frac{\tau(r_1)}{\tau(r_1)+[1-\tau(r_1)]\eu}<\pi<\frac{\tau(r_2)}{\tau(r_2)+[1-\tau(r_2)]\eu},
  \]
  with the endpoints of Chapter~\ref{app:numbers}. The high-cost and trading-cost windows alone hold exactly
  between the same lower endpoint and \(\tau(r_0)\).
\item At the example's slot (\(1+u=\eu\) and the example's \(\beta\)) and \(\pi=\tfrac12\), entry, high-class
  ownership and class \(H\)'s share of the acquisitions won by the challenger take, at the three strengths, the
  values of Chapter~\ref{app:numbers}.
\item For these point values and the strength range \([r_0,r_2]\), decisive advantage on the full uniform family
  and the affordability condition \(h-r_2\ge c_H\) of Proposition~3(c) hold together exactly on the cost window
  \[
  \frac{h^2-p^2}{2h}\le c_H\le h-r_2 .
  \]
  The example's cost lies in this window, with positive margins. For these class values, reserve and cost, every
  economy with a uniform incumbent \(U[0,r]\) in which some order recruits a high-cost challenger has
  \[
  r<h-c_H+\sqrt{(h-c_H)^2-p^2}<h,
  \]
  whatever the slot, prior, \(\rho\), \(k\) and research.
\end{enumerate}
\end{paperresult}

The short form in the main text (Section~4.2 of the paper) adds one line: on the band, class \(H\) outbids every
incumbent on the path, entry rises from \(\rho\) at \(r_0\) to \(\rho+(1-\rho)(b+\pi\lambda)\) at \(r_1\) and returns
to \(\rho\) at \(r_2\), and high-class ownership rises from \(\pi\rho\) to \(\pi[\rho+(1-\rho)(b+\lambda)]\).

In plain words: at the prior \(\tfrac12\), Proposition~2 applies in the example exactly when the top posterior falls
in a band whose ends are the thresholds of the strong and the very strong incumbent (a). Outside the band one
condition fails, and at the lower edge the tie rule matters (b). With the example's slot, the band is a range of
priors (c). Part (e) makes class \(H\) outbid every incumbent against which a price recruits
(Section~\ref{sec:example-decisive}).

\subsection{Why the conditions reduce to a band}\label{sec:example-reduction}

The proof of part (a) (Appendix~A.3 of the paper) checks each of the five inequalities in (H1)--(H3). Two of them do
not involve the slot. The other three each give one bound on \(u\).

\paragraph{Step 1: the posterior bounds at the prior one half.} Put \(\pi=\tfrac12\) in \eqref{eq:beliefs-bounds}.
\begin{align*}
M(\tfrac12)&=\frac{\tfrac12(1+u)}{\tfrac12(1+u)+\tfrac12} \why{\eqref{eq:beliefs-bounds}}\\
&=\frac{1+u}{2+u} \why{multiply above and below by 2}\\[4pt]
m(\tfrac12)&=\frac{\tfrac12}{\tfrac12+\tfrac12(1+u)}=\frac{1}{2+u} \why{same steps}\\[4pt]
1-M(\tfrac12)&=\frac{2+u-1-u}{2+u}=\frac{1}{2+u} \why{common denominator}
\end{align*}
So \(m(\tfrac12)=1-M(\tfrac12)\), and \(\kappa(\tfrac12)=1/(2+u)\). Also \(M(\tfrac12)=1-\tfrac1{2+u}\) rises strictly
with \(u\).

\paragraph{Step 2: the two conditions without the slot.} The left side of (H1), \(B_{r_0}(\tfrac12)<c_H\), contains
only the auction at \(r_0\), the prior and the high cost. The left side of (H2), \(\DT(r_0)<k\), contains only the
spread at \(r_0\) and the trading cost. Neither depends on \(\lambda\) or \(\beta\). Both hold at the example's
values, with positive margins (Appendix~A.3 of the paper). The first one says \(\tau(r_0)>\tfrac12\).

\paragraph{Step 3: the right side of (H1).} At \(r_1\), \(g_H>g_L\), so \(B_{r_1}\) rises strictly with the belief,
and \(c_H<B_{r_1}(M)\) holds exactly when \(M>\tau(r_1)\) (Section~\ref{sec:beliefs-threshold}). Turn this into a
bound on \(u\). Both thresholds on the path lie below one at the example's values, so \(1-\tau>0\).
\begin{align*}
M(\tfrac12)>\tau&\iff\frac{1+u}{2+u}>\tau \why{Step 1}\\
&\iff 1+u>\tau(2+u) \why{multiply by \(2+u>0\)}\\
&\iff u-\tau u>2\tau-1 \why{subtract \(1+\tau u\)}\\
&\iff u>\frac{2\tau-1}{1-\tau} \why{divide by \(1-\tau>0\)}
\end{align*}
With \(\tau=\tau(r_1)\), the right side of (H1) is \(u>u_{\rm lo}\).

\paragraph{Step 4: the exclusion condition.} The same steps with the opposite inequality give
\(B_{r_2}(M)<c_H\iff M<\tau(r_2)\iff u<u_{\rm hi}\).

\paragraph{Step 5: the right side of (H2).}
\begin{align*}
k<\rho\,\kappa(\tfrac12)\,\DT(r_1)&\iff k<\frac{\rho\,\DT(r_1)}{2+u} \why{Step 1}\\
&\iff 2+u<\frac{\rho\,\DT(r_1)}{k} \why{multiply by \((2+u)/k>0\)}\\
&\iff u<\frac{\rho\,\DT(r_1)}{k}-2=u_k \why{subtract 2}
\end{align*}
At the example's values, \(u_{\rm hi}<u_k\). So Step 4 implies Step 5, and the right side of (H2) adds no further
restriction inside the band.

\paragraph{Step 6: combine.} Given Step 2, the five inequalities hold together exactly when
\(u_{\rm lo}<u<u_{\rm hi}\). Steps 3 and 4 show that this is the same as \(\tau(r_1)<M(\tfrac12)<\tau(r_2)\): the top
posterior lies between the thresholds of the strong and the very strong incumbent.

\paragraph{Step 7: an interval of presence probabilities.} From \(u=\lambda/[(1-\lambda)\beta]\), solve for
\(\lambda\) at a fixed \(\beta\).
\begin{align*}
u\beta(1-\lambda)&=\lambda \why{multiply by \((1-\lambda)\beta\)}\\
u\beta&=\lambda(1+u\beta) \why{add \(u\beta\lambda\) to both sides}\\
\lambda&=\frac{u\beta}{1+u\beta}=1-\frac{1}{1+u\beta} \why{divide by \(1+u\beta\)}
\end{align*}
At a fixed \(\beta\), \(\lambda\) rises strictly with \(u\). So the band of \(u\) is an open interval of
\(\lambda\).

\begin{keyidea}
At the prior \(\tfrac12\), the conditions of Proposition~2 contain the slot only through the top posterior
\(M(\tfrac12)=(1+u)/(2+u)\) and the factor \(\kappa(\tfrac12)=1/(2+u)\). Both are functions of \(u\) alone. Two slots
with the same \(u\) and different \((\lambda,\beta)\) give the same answer.
\end{keyidea}

The slot enters the outcomes in more ways. Entry contains \(b\) and \(\lambda\) separately. Use \eqref{eq:ccc-entry}
at \(\pi=\tfrac12\):
\begin{align*}
\Ent_1&=\rho+(1-\rho)(b+\pi\lambda) \why{\eqref{eq:ccc-entry}}\\
&=\rho+(1-\rho)\Bigl(b+\frac\lambda2\Bigr) \why{\(\pi=\tfrac12\)}
\end{align*}
The term \(b+\lambda/2\) is the probability of a purchase: a liquidity purchase with probability \(b\) in both
classes, plus an informed purchase with probability \(\lambda\) in class \(H\), which has prior weight one half.
Entry includes the floor \(\rho\) of the zero-cost challenger. At \(r_0\) and \(r_2\), entry is \(\rho\)
(Proposition~2(i) and (iii)).

For ownership, recall \(\Own_H=\pi\,e_H\,W_H(F)\), with \(e_H\) the entry in class \(H\) (Chapter~\ref{ch:model}).
Every strength on the path lies below \(h\), so \(W_H=1\) at all three. Proposition~2(ii) gives
\(e_H=\rho+(1-\rho)(b+\lambda)\) at \(r_1\), and \(e_H=\rho\) at \(r_0\) and \(r_2\). So ownership is
\(\rho/2\), \([\rho+(1-\rho)(b+\lambda)]/2\) and \(\rho/2\).

For part (d), the share of class \(H\) among the acquisitions won by the challenger is \(\Own_H/(\Own_H+\Own_L)\),
with \(\Own_L=(1-\pi)\,e_L\,W_L(F)\). Class \(L\) wins when \(R<\ell\), so \(W_L(F_r)=F_r(\ell)=\ell/r\). The share is
\[
\frac{\pi\,e_H}{\pi\,e_H+(1-\pi)\,e_L\,\ell/r}
\]
(Appendix~A.3 of the paper). From \(r_0\) to \(r_1\), the ratio \(e_H/e_L\) rises and \(\ell/r\) falls, so the share
rises. Proposition~2(ii) says the same: the share rises strictly when \(W_L(F_0)>0\), and here \(W_L(F_0)=\ell/r_0>0\).

\subsection{Outside the band and other priors}

Part (b) reads each edge. Below the band, orders are still full at \(r_1\), but \(M<\tau(r_1)\), so no order
recruits and entry at \(r_1\) stays at \(\rho\); this is the case that Chapter~\ref{ch:ccc} mentions when the right
side of (H1) fails. At the lower edge, \(M=\tau(r_1)\), and a purchase leaves the high-cost challenger indifferent:
the tie rule decides. Above the band, \(M\ge\tau(r_2)\), so a purchase at \(r_2\) recruits under the default tie
rule. For \(u\ge u_k\), the trading bound fails to force full orders at \(r_1\).

Part (c) fixes the slot at \(1+u=\eu\) and lets the prior vary. Recall that the posterior rises with its odds
(Section~\ref{sec:toolkit-odds}), and that odds \(o\) give the probability \(o/(1+o)\).

\begin{derivation}[the top posterior above a threshold, as a condition on the prior]
\begin{align*}
M(\pi)>\tau&\iff\frac{M(\pi)}{1-M(\pi)}>\frac{\tau}{1-\tau} \why{odds rise with the probability}\\
&\iff\frac{\pi(1+u)}{1-\pi}>\frac{\tau}{1-\tau} \why{the odds of \(M(\pi)\)}\\
&\iff\frac{\pi}{1-\pi}>\frac{\tau}{(1-\tau)(1+u)} \why{divide by \(1+u>0\)}\\
&\iff\pi>\frac{\tau}{\tau+(1-\tau)(1+u)} \why{odds \(o\) give \(o/(1+o)\)}
\end{align*}
\end{derivation}

With \(\tau=\tau(r_1)\) and \(1+u=\eu\), the right side of (H1) gives the lower end of the range in part (c). With
\(\tau=\tau(r_2)\) and the reverse inequality, (H3) gives the upper end. The left side of (H1) is
\(B_{r_0}(\pi)<c_H\), that is \(\pi<\tau(r_0)\). The left side of (H2) does not contain \(\pi\). The right side of
(H2) gives a range of priors on which \(\kappa(\pi)\) is large enough, and at the example's values this range
contains the other two (Appendix~A.3 of the paper). So the range of part (c) remains. Without (H3) its upper end
is \(\tau(r_0)\). The range of priors in part (c) is narrow. Its upper end lies just above the example's prior.

\subsection{An open set of parameters}\label{sec:example-open}

The paper adds that (H1)--(H3) hold on an open set of parameters that contains the example (Section~4.2 and
Appendix~A.3 of the paper). The argument is the open-set fact of Section~\ref{sec:toolkit-monotone}: a finite list of
strict inequalities between continuous functions survives small changes of the parameters.

Keep the example's classes and reserve, with each incumbent uniform on \([0,r_i]\). Each side of (H1)--(H3)
is a continuous function of \((r_0,r_1,r_2,c_H,k,\rho,\pi,\lambda,\beta)\): \(B_r(\mu)\) and \(\DT(r)\) by
\eqref{eq:example-payoffs}; \(u=\lambda/[(1-\lambda)\beta]\), and with it \(M(\pi)\) and \(m(\pi)\); and
\(\kappa\), a minimum of two continuous functions. All five inequalities are strict at the example. So every
parameter list near enough to the example keeps all five, and Proposition~2 still applies. The set is open in the
domain \(0<\beta<1/2\). In words, the reversal in the example is not a knife-edge case. The fact says nothing about
the size of the set, and the band of Chapter~\ref{app:numbers} is narrow.

\begin{checkyourself}
\label{q:example-ceiling}
Take a slot in the band, at the prior \(\tfrac12\). The entry ceiling \(r_C\) is the strength at which
\(\tau(r)=M(\tfrac12)\) (Section~\ref{sec:example-ceiling}).
(a)~Show that \(r_C\) lies strictly between \(r_1\) and \(r_2\).
(b)~Explain why a larger \(u\) inside the band moves \(r_C\) to a higher strength.
(c)~Say what happens to entry at \(r_2\) when \(u\) passes \(u_{\rm hi}\), with \(u<u_k\) kept.
\end{checkyourself}

% ====================================================================
\section{Decisive advantage in the example}\label{sec:example-decisive}

In the example's three economies class \(H\) outbids the incumbent, because \(r_2<h\). Decisive advantage asks
more: does class \(H\) outbid \emph{every} uniform incumbent against which a price recruits? Recall from Lemma~A.6
(Section~\ref{sec:ccc-decisive}) that a family of incumbent laws has decisive advantage when \(g_H(F)>c_H\) implies
\(W_H(F)=1\) for every law \(F\) in it, and that the full uniform family has it exactly when
\(c_H\ge(h^2-p^2)/(2h)\). Here is the short version in the example's notation.

\begin{derivation}[decisive advantage for the uniform family]
\emph{The chance to win.} For \(r\le h\), \(W_H(F_r)=F_r(h)=1\). For \(r>h\), \(W_H(F_r)=h/r<1\).

\emph{Gross profit of class \(H\) at every strength.} By \eqref{eq:auction-gross}, \(g_H(r)=\int_p^hF_r(y)\,dy\).
\begin{align*}
r\le p:\quad g_H&=\int_p^h1\,dy=h-p \why{\(F_r=1\) above \(r\)}\\
p<r\le h:\quad g_H&=h-\frac r2-\frac{p^2}{2r} \why{\eqref{eq:example-payoffs}}\\
r\ge h:\quad g_H&=\int_p^h\frac yr\,dy=\frac{h^2-p^2}{2r} \why{\(F_r(y)=y/r\) on \([p,h]\)}
\end{align*}
The middle line at \(r=h\) gives \(h-\tfrac h2-\tfrac{p^2}{2h}=\tfrac{h^2-p^2}{2h}\), the value of the last line at
\(r=h\). So \(g_H\) is continuous. It falls strictly for \(r>p\): the middle line has \(g_H'<0\), and the last line
falls like \(1/r\).

\emph{The equivalence.} Write \(\bar g=g_H(h)=(h^2-p^2)/(2h)\). If \(c_H\ge\bar g\), every \(r\) with
\(g_H(r)>c_H\ge g_H(h)\) has \(r<h\), since \(g_H\) falls strictly above \(p\); there \(W_H=1\). If \(c_H<\bar g\),
continuity gives \(g_H(r)>c_H\) at strengths a little above \(h\), where \(W_H=h/r<1\).
\end{derivation}

\begin{intuition}[what the cost floor means]
The number \((h^2-p^2)/(2h)\) is class \(H\)'s gross profit against the uniform incumbent on \([0,h]\), the strongest
uniform incumbent that class \(H\) still beats for sure. If the high cost is at least this profit, then gross profit
can exceed the cost only against weaker incumbents, and class \(H\) beats all of them.
\end{intuition}

\emph{What it buys.} A recruited high-cost challenger needs \(B(\bar\mu)\ge c_H\) at an interior belief
\(\bar\mu<1\), and Lemma~A.6 then gives \(g_H>c_H\) (Section~\ref{sec:ccc-decisive}). The example's high cost exceeds
\((h^2-p^2)/(2h)\) (Section~4.2 of the paper). So, against every uniform incumbent at which a price recruits a
high-cost challenger, class \(H\) outbids the incumbent. High-class ownership is then \(\pi e_H\), and the
takeover-share conclusions of Proposition~2(ii) apply: the strong incumbent takes wins only from the low class.

\emph{The recruitment bound of part (e).} Recruitment needs \(g_H(r)>c_H\). Find where \(g_H\) crosses \(c_H\). For
\(r\le p\), \(g_H=h-p\), and \(h-p>c_H\) at the example's values (Appendix~A.3 of the paper). For \(r\ge h\),
\(g_H\le\bar g<c_H\), since the example's cost exceeds \(\bar g\). So \(g_H\) meets \(c_H\) in \((p,h)\), where the middle formula applies.
\begin{align*}
h-\frac r2-\frac{p^2}{2r}=c_H&\iff 2rh-r^2-p^2=2rc_H \why{multiply by \(2r>0\)}\\
&\iff r^2-2(h-c_H)\,r+p^2=0 \why{rearrange}\\
&\iff r=(h-c_H)\pm\sqrt{(h-c_H)^2-p^2} \why{quadratic formula}
\end{align*}
The roots are real because \(h-c_H>p\). Their product is \(p^2\), so the lower root lies below \(p\) and the upper
root lies above \(p\). The upper root is the only point where \(g_H=c_H\), and since \(g_H\) falls strictly, \(g_H(r)>c_H\) exactly
when \(r<h-c_H+\sqrt{(h-c_H)^2-p^2}\). This bound lies below \(h\):
\begin{align*}
h-c_H+\sqrt{(h-c_H)^2-p^2}<h&\iff\sqrt{(h-c_H)^2-p^2}<c_H \why{subtract \(h-c_H\)}\\
&\iff h^2-2hc_H-p^2<0 \why{square; both sides \(\ge0\)}\\
&\iff c_H>\frac{h^2-p^2}{2h} \why{divide by \(2h>0\)}
\end{align*}
The example's cost satisfies the last line strictly. The bound uses only the auction and \(c_H\), so it holds
whatever the slot, prior, \(\rho\), \(k\) and research.

The upper end of the cost window, \(c_H\le h-r_2\), is the affordability condition of Proposition~3(c).
Chapter~\ref{ch:coexistence} uses it.

% ====================================================================
\section{The entry ceiling and the width of the band}\label{sec:example-ceiling}

\subsection{The entry ceiling}

The entry ceiling \(r_C\) is the strength at which the threshold meets the top posterior, \(\tau(r_C)=M(\tfrac12)\)
(Section~4.2 of the paper). On the band, \(\tau(r_1)<M<\tau(r_2)\). The threshold is continuous and rises strictly,
so the intermediate value theorem (Section~\ref{sec:toolkit-monotone}) gives exactly one such strength, and it lies
between \(r_1\) and \(r_2\). In the example, \(r_C\) lies just below \(r_2\).

\emph{The exclusion condition holds above the ceiling.} Gross profit at every belief falls strictly with strength
above \(\ell\) (Section~\ref{sec:example-shape}). At the ceiling, \(B_{r_C}(M)=c_H\). So for every \(r\) in
\((r_C,h)\),
\begin{align*}
B_r(M)&<B_{r_C}(M) \why{\(B\) falls strictly in \(r\); \(r>r_C\)}\\
&=c_H \why{\(\tau(r_C)=M\)}
\end{align*}
So (H3) holds with any \(r_2\) in \((r_C,h)\) (Section~4.2 of the paper). Parts (i) and (ii) of Proposition~2 do
not use the exclusion condition, so they hold whatever \(r_2\) is.

\begin{figure}[htbp]
\centering
\includegraphics[width=\linewidth]{beliefs}
\caption{Beliefs and the entry ceiling in the example, at the prior \(\tfrac12\). Panel (a): the preparation
threshold \(\tau(r)\) rises with strength from \(r_0\) to \(r_2\); the grey band holds every posterior between
\(m(\tfrac12)\) and \(M(\tfrac12)\); the threshold meets the top of the band at the entry ceiling \(r_C\). Panel (b):
the strengths from \(r_1\) to \(r_2\) on a finer vertical scale; the shaded band holds every top posterior that puts
\(r_C\) between \(r_1\) and \(r_2\). The printed values are those of Chapter~\ref{app:numbers}.}
\label{fig:example-beliefs}
\end{figure}

\subsection{How to read the figure}

Figure~\ref{fig:example-beliefs} is the paper's Figure~1. In panel (a), the dotted line at the prior \(\tfrac12\)
lies below the threshold at every strength, so an uninformative price never recruits a high-cost challenger on the
path. The grey band holds every posterior that an investor strategy can produce, when \(\beta\le1/3\) (Lemma~A.1).
The threshold curve rises until it meets the top of the band at \(r_C\), just below \(r_2\). To the left of
\(r_C\), a purchase that lifts the belief to \(M\) recruits; to the right, nothing recruits. Panel (b) enlarges the
strengths from \(r_1\) to \(r_2\) on a much finer vertical scale. Its shaded band runs from \(\tau(r_1)\) to
\(\tau(r_2)\), the band of Corollary~1(a), and the example's top posterior lies just below its upper edge.

\subsection{Why the band is narrow}

Both ends of the band are thresholds, so the width of the band is the rise of the threshold from \(r_1\) to
\(r_2\). The paper explains the narrow band with one sentence: the threshold rises slowly, because class \(H\)'s expected payment
\(r/2+p^2/(2r)\) rises by less than one half per unit of strength, a change that is small next to the gap
\(g_H-g_L\) (Section~4.2 of the paper). The next derivation is the guide's own form of this sentence.

\begin{derivation}[the slope of the threshold]
The threshold moves so that gross profit at the threshold stays equal to the cost: \(B_r(\tau(r))=c_H\) at every
\(r\). Differentiate both sides in \(r\). The left side depends on \(r\) directly and through \(\tau(r)\).
\begin{align*}
0&=\frac{d}{dr}\Bigl[g_L(r)+\tau(r)\,\gamma(r)\Bigr] \why{\(B_r(\tau)=g_L+\tau\gamma\); \(c_H\) fixed}\\
&=g_L'+\tau'\gamma+\tau\gamma' \why{product rule}\\
&=g_L'+\tau(g_H'-g_L')+\tau'\gamma \why{\(\gamma'=g_H'-g_L'\)}\\
&=(1-\tau)\,g_L'+\tau\,g_H'+\tau'\gamma \why{collect \(g_L'\)}
\end{align*}
Solve for \(\tau'\):
\begin{equation*}
\tau'(r)=\frac{(1-\tau)\,\bigl(-g_L'(r)\bigr)+\tau\,\bigl(-g_H'(r)\bigr)}{\gamma(r)} .
\end{equation*}
The numerator is the fall of gross profit at the belief \(\tau\), per unit of strength.
\end{derivation}

Next, bound the two falls. From Section~\ref{sec:example-shape}:
\begin{align*}
-g_H'(r)&=t_H'(r)=\frac12-\frac{p^2}{2r^2}\in\Bigl(0,\frac12\Bigr) \why{class \(H\)'s payment}\\
-g_L'(r)&=\frac{\ell^2-p^2}{2r^2}<\frac{\ell^2}{2r^2}<\frac12 \why{\(r>\ell\)}
\end{align*}
and \(-g_L'(r)>0\). Where \(0\le\tau\le1\), the numerator of \(\tau'\) is an average of two numbers in
\((0,\tfrac12)\), with weights \(1-\tau\) and \(\tau\). So it lies in \((0,\tfrac12)\), and
\[
0<\tau'(r)<\frac{1}{2\gamma(r)} .
\]
In words: per unit of strength, the threshold rises by less than one half divided by the profit gap. The gap
\(\gamma=g_H-g_L=h-\ell-\DT\) is large when the class values lie far apart. Then the threshold barely moves, and two
strengths a modest distance apart give two thresholds close together. The band of top posteriors between them is
narrow.

At the other extreme, a small profit gap would let a small fall in profit move the threshold a lot, and the band
would be wide.

\begin{checkyourself}
\label{q:example-width}
Use \(0<\tau'(r)<1/[2\gamma(r)]\) on \([r_1,r_2]\), and the fact that \(\gamma\) falls with \(r\), to show that
\[
\tau(r_2)-\tau(r_1)<\frac{r_2-r_1}{2\gamma(r_2)} .
\]
Say in words why a large profit gap makes the band of Corollary~1 narrow.
\end{checkyourself}

% ====================================================================
\section{The fixed-information control}\label{sec:example-control}

The second part of the chapter asks where the rise in entry comes from. Proposition~2 compares economies in which
two things change together: the incumbent is stronger, and the information in the price is different. A control
holds one of them fixed.

\begin{intuition}[hold the information fixed]
Suppose the challenger learns the same thing about its class whatever the incumbent. A stronger incumbent lowers
gross profit at every belief (Proposition~1(a)), so a challenger that prepares against the strong incumbent also
prepares against the weak one. Entry can only fall.
\end{intuition}

\begin{analogy}[the loudspeaker again]
Chapter~\ref{ch:ccc} pictured the informed investor as a loudspeaker. Panel~B of the paper's outcome table fixes the
message at every strength. Panel~C keeps the loudspeaker on but plugs the listener's ears.

Where the analogy breaks: at the weak incumbent the investor would not trade, because trade does not pay there.
Panel~B forces the message anyway, so at that strength it is not an equilibrium.
\end{analogy}

\subsection{The statement}

\needspace{10\baselineskip}
\begin{paperresult}{Proposition A.2 (Fixed information)}
Fix the joint law of the class \(j\) and a signal \(S\) that the challenger observes. Fix the class value laws
\(G_j\) and the cost law. In each economy \(i\), require that, conditional on \((j,S)\), the joint law of
\((\theta,R)\) is \(G_j\otimes F_i\). The cost is independent of the class, both values and \(S\). If the incumbent
strengthens from \(F_0\) to \(F_1\), entry is nonincreasing when an indifferent high-cost challenger prepares. Its
decrease is exactly
\[
\Ent_0-\Ent_1=\Pr\bigl\{B_{F_1}(\mu(S))<C\le B_{F_0}(\mu(S))\bigr\},\qquad \mu(S)=\Pr(j=H\mid S),
\]
and the decrease is strict exactly when this probability is positive.
\end{paperresult}

In plain words, line by line:
\begin{itemize}
\item \emph{The signal.} \(S\) is whatever the challenger sees before it decides. Its joint law with the class is
  the same in both economies. So the challenger's belief \(\mu(S)\) has the same law in both economies.
\item \emph{The product law \(G_j\otimes F_i\).} Given the class and the signal, the challenger's value has the law
  \(G_j\), the incumbent's value has the law \(F_i\), and the two are independent (Section~4.2 of the paper).
\item \emph{The cost} has any law, the same in both economies, and is independent of everything else.
\item \emph{The conclusion.} Under the default tie rule, entry cannot rise. It falls by the probability that the
  cost lies above gross profit against the strong incumbent and at or below gross profit against the weak one.
\end{itemize}

\subsection{The proof}

\paragraph{Step 1: the belief at the decision.} The cost is independent of the class, the values and \(S\). So when
the challenger learns its cost, it learns nothing more about its class. Its complete class belief at the decision is
\(\mu(S)=\Pr(j=H\mid S)\).

\paragraph{Step 2: gross profit at the decision.} A prepared challenger keeps \((\theta-\max\{p,R\})^+\); for
\(\theta<p\) this is zero, as for a challenger that does not bid. Average over the class and then over the values.
\begin{align*}
\E\bigl[(\theta-\max\{p,R\})^+\mid S\bigr]
&=\sum_{j}\Pr(j\mid S)\,\E\bigl[(\theta-\max\{p,R\})^+\mid j,S\bigr] \why{total expectation}\\
&=\bigl(1-\mu(S)\bigr)\,g_L(F_i)+\mu(S)\,g_H(F_i) \why{given \((j,S)\): \(G_j\otimes F_i\)}\\
&=B_{F_i}\bigl(\mu(S)\bigr) \why{definition of \(B\)}
\end{align*}
The second line uses the product law: given \((j,S)\), the pair \((\theta,R)\) has exactly the law that defines
\(g_j(F_i)\) (Chapter~\ref{ch:auction}).

\paragraph{Step 3: the preparation rule and entry.} The challenger compares its cost with gross profit. When an
indifferent high-cost challenger prepares, it prepares exactly when \(C\le B_{F_i}(\mu(S))\). Entry is the probability
of this event:
\[
\Ent_i=\Pr\bigl\{C\le B_{F_i}(\mu(S))\bigr\}=\E\bigl[\1\{C\le B_{F_i}(\mu(S))\}\bigr].
\]

\paragraph{Step 4: gross profit is lower against the stronger incumbent.} Proposition~1(a)
(Section~\ref{sec:auction-prop1}) gives \(B_{F_1}(\mu)\le B_{F_0}(\mu)\) at every belief \(\mu\). Write
\(B_0=B_{F_0}(\mu(S))\) and \(B_1=B_{F_1}(\mu(S))\) for short, so \(B_1\le B_0\).

\paragraph{Step 5: a pointwise identity.} Compare the two indicators at each value of \(C\). If \(C\le B_1\), then
also \(C\le B_0\), and both are one. If \(B_1<C\le B_0\), the first is one and the second is zero. If \(C>B_0\), both
are zero. So, at every outcome,
\[
\1\{C\le B_0\}-\1\{C\le B_1\}=\1\{B_1<C\le B_0\}\ge0 .
\]

\paragraph{Step 6: take expectations.} The joint law of \((S,C)\) is the same in both economies: the law of \((j,S)\)
is fixed, and \(C\) is independent of \(S\) with a fixed law. The function \(\mu(\cdot)\) is the same, too. Only the
function \(B_{F_i}\) changes.
\begin{align*}
\Ent_0-\Ent_1&=\E\bigl[\1\{C\le B_0\}\bigr]-\E\bigl[\1\{C\le B_1\}\bigr] \why{Step 3}\\
&=\E\bigl[\1\{C\le B_0\}-\1\{C\le B_1\}\bigr] \why{linearity; same law of \((S,C)\)}\\
&=\E\bigl[\1\{B_1<C\le B_0\}\bigr] \why{Step 5}\\
&=\Pr\bigl\{B_{F_1}(\mu(S))<C\le B_{F_0}(\mu(S))\bigr\}\ge0 \why{expectation of an indicator}
\end{align*}
Entry does not rise, and it falls strictly exactly when this probability is positive.

\paragraph{Step 7: another tie rule.} Suppose an indifferent high-cost challenger prepares with a probability \(q\)
that is the same in both economies, so it prepares with probability \(\1\{C<B_i\}+q\,\1\{C=B_i\}\). With
\(B_1\le B_0\), the difference of the two expressions is \(0\) for \(C<B_1\) and for \(C>B_0\), \(1-q\) at
\(C=B_1<B_0\), \(1\) strictly between, \(q\) at \(C=B_0>B_1\), and \(0\) when the two profits are equal. Each value
is nonnegative, so entry does not rise under this common tie probability either (Appendix~A.6 of the paper).

\emph{The two-point cost.} With \(C\in\{0,c_H\}\), the cost \(C=0\) never lies in the interval
\((B_1,B_0]\), because gross profit is never negative, \(B_1\ge0\). The cost is \(c_H\) with probability
\(1-\rho\), independently of \(S\). So
\[
\Ent_0-\Ent_1=(1-\rho)\,\Pr\bigl\{B_{F_1}(\mu(S))<c_H\le B_{F_0}(\mu(S))\bigr\}.
\]
In words: entry falls by the high-cost share times the probability of a signal whose belief recruits against the
weak incumbent but not against the strong one.

\begin{watchout}
The control keeps three independence conditions, and each one is needed.
\begin{itemize}
\item \emph{The signal says nothing about the values beyond the class.} If \(S\) carried news about \(\theta\) given
  the class, gross profit at the decision would not be \(B_{F_i}(\mu(S))\), and Step 2 would fail. The same
  condition appears in the paper's signal extension (Chapter~\ref{ch:robustness}).
\item \emph{The incumbent's value is independent of the class and the signal.} Otherwise the signal could carry news about
  \(R\), and gross profit at the decision would not be computed with \(F_i\).
\item \emph{The cost is independent of everything.} Otherwise the cost would carry news about the class, and
  \(\mu(S)\) would not be the belief at the decision (Step 1).
\end{itemize}
\end{watchout}

\subsection{The control in the order slot}\label{sec:example-slot}

In the model the challenger sees a price, and the price comes from the order slot. The slot version of the
experiment fixes \(\lambda\), \(\beta\), \(\pi\), the cost law and both order distributions (Section~4.2 of the
paper). The appendix uses full orders, \((\sigma_H,\sigma_L)=(\delta_1,\delta_{-1})\): class \(H\) buys and class
\(L\) sells for sure (Appendix~A.6 of the paper). The question is whether the price then gives the challenger the
same information at every strength.

\begin{derivation}[competitive prices reveal the same posterior at every strength]
\emph{Step 1: posteriors.} With full orders, Lemma~A.1(d) gives the posteriors \((m,\pi,M)\) after the orders
\((-1,0,1)\). They depend only on \(\lambda\), \(\beta\), \(\pi\) and the orders, not on \(r\).

\emph{Step 2: entry and prices.} Let entry follow the threshold rule under the default tie rule,
\(e_r(\mu)=\rho+(1-\rho)\,\1\{B_r(\mu)\ge c_H\}\). Since \(g_H\ge g_L\), \(B_r\) never falls with the belief, so a
higher posterior never gets lower entry. Set the competitive price \eqref{eq:beliefs-price},
\(P(x)=t_0+e_r(\mu(x))\,w(\mu(x))\), with \(w(\mu)=t_L-t_0+\DT\mu\).

\emph{Step 3: each price reveals its posterior.} The paragraph on constructed tuples
(Section~\ref{sec:beliefs-constructed}) applies when \(\DT>0\): entry is in \([\rho,1]\) and never lower at a higher
posterior. It gives prices that rise strictly with the posterior. So the three prices differ, and each price reveals
its posterior.

\emph{Step 4: the signal.} The challenger's belief is \(\mu(X)\), a function of the order. The law of the order given
the class does not depend on \(r\). Given the class and the order, the values have the law \(G_j\otimes F_r\),
because the order depends only on the class and the slot draws. So the hypotheses of Proposition~A.2 hold, with the
order as the signal.
\end{derivation}

\begin{watchout}
The slot version keeps the price information fixed only if the spread is positive in each economy that it compares
(Appendix~A.2 and A.6 of the paper; Section~\ref{sec:beliefs-reveal}).
\begin{itemize}
\item With \(\DT=0\), the price is \(P(x)=t_0+e(x)(t_L-t_0)\). It does not move with the posterior. Two orders with
  different posteriors and the same entry would share one price, and the challenger would see only an average of
  their posteriors. The information would then differ from the economies with \(\DT>0\).
\item The entry floor \(\rho>0\) keeps the posterior term \(e(x)\DT\mu(x)\) in the price.
\item In the example, \(\DT(r)=(r-\ell)^2/(2r)>0\) at every strength above \(\ell\). The spread at the weak
  incumbent is small, but it is positive, so the control works there too.
\end{itemize}
\end{watchout}

\subsection{The control in the example}

Apply the control to the example at the prior \(\tfrac12\), with full orders at every strength on the path.

\emph{Only the top posterior can recruit.} At the example's values, the threshold at \(r=\ell\) exceeds the prior,
\[
\tau(\ell)=\frac{2\ell c_H-(\ell^2-p^2)}{2\ell(h-\ell)}>\pi
\]
(Appendix~A.6 of the paper), and \(\tau\) rises on \([\ell,h]\). So \(\tau(r)>\pi>m\) at every strength. After a sale or
no order, entry is \(\rho\). After a purchase, the belief is \(M\), and it recruits exactly when \(M\ge\tau(r)\), that
is, when \(r\le r_C\).

\emph{Entry with information held fixed.} A purchase has probability \(\pi(b+\lambda)+(1-\pi)b=b+\pi\lambda\).
\begin{align*}
r\le r_C:\quad\Ent&=(b+\pi\lambda)\cdot1+\bigl[1-(b+\pi\lambda)\bigr]\rho \why{a purchase recruits}\\
&=\rho+(1-\rho)\Bigl(b+\frac\lambda2\Bigr) \why{rearrange; \(\pi=\tfrac12\)}\\
r>r_C:\quad\Ent&=\rho \why{no order recruits}
\end{align*}
So entry is \(\rho+(1-\rho)(b+\lambda/2)\) on \((\ell,r_C]\) and \(\rho\) above \(r_C\).

\emph{The probability of Proposition~A.2.} With the two-point cost, the fall is \((1-\rho)\) times the probability
of an order whose posterior recruits against the weaker incumbent and not against the stronger one.
\begin{itemize}
\item From \(r_0\) to \(r_1\): the sale and no order recruit at neither strength, and a purchase recruits at both,
  since \(r_0<r_1<r_C\). No order switches. The probability is zero, and entry stays the same.
\item From \(r_1\) to \(r_2\): a purchase recruits at \(r_1\le r_C\) and not at \(r_2>r_C\). The probability is
  \((1-\rho)\Pr(X=1)=(1-\rho)(b+\lambda/2)\), and entry falls by exactly this amount, to \(\rho\).
\end{itemize}

\begin{watchout}
Panel~B imposes full orders. It is a control, and at the weak incumbent it is not an equilibrium. There, the
purchase earns \(\Delta_0(1-M)-k\) and the sale earns \(\rho\Delta_0m-k\), and both are negative because
\(\Delta_0<k\). So the investor would rather abstain (Appendix~A.6 of the paper). At \(r_1\), the bound
\(\rho\kappa\Delta_1>k\) forces full orders in every candidate (Step~2 of the proof of Proposition~2), so the
imposed profile is the unique equilibrium there. Keep the distinction of Chapter~\ref{ch:model}: the control fixes
the orders across economies, while each equilibrium comparison solves every economy again.
\end{watchout}

\begin{checkyourself}
\label{q:example-control}
In the slot version of the control in the example, at the prior \(\tfrac12\), take two strengths \(r<r'\) in
\((\ell,h)\). Find the probability of Proposition~A.2 in three cases: (a)~\(r<r'\le r_C\); (b)~\(r\le r_C<r'\);
(c)~\(r_C<r<r'\). In each case, say in words why entry changes or does not change.
\end{checkyourself}

% ====================================================================
\section{The two panels and what they show}\label{sec:example-panels}

The paper's outcome table (Table~3 of the paper) has three panels for the example. Panel~A lists the equilibrium of
Proposition~2 at \(r_0\), \(r_1\) and \(r_2\): entry rises from \(\rho\) and falls back to \(\rho\). Panels~B and~C are
the two controls. Both compare \(r_0\) with \(r_1\). The numbers are in Chapter~\ref{app:numbers}; this section reads
the panels in words.

\subsection{Panel B: full orders held fixed}

Panel~B imposes full orders at both strengths, so the posteriors are the same at \(r_0\) and \(r_1\), and the
prices reveal them. Only the top posterior can recruit, and it recruits at both strengths, because both lie below
the entry ceiling. So entry is the same at the two strengths, and it equals the entry of Panel~A at \(r_1\). With information fixed,
entry stays constant from \(r_0\) to \(r_1\) and falls only above the entry ceiling \(r_C\). This is
Proposition~A.2 at work: at fixed information, a stronger incumbent never raises entry.

\subsection{Panel C: the challenger does not see the price}

Panel~C hides the price from the challenger. The challenger then decides on its prior. Gross profit at the prior
falls short of the high cost at both strengths: \(B_{r_1}(\pi)\le B_{r_0}(\pi)<c_H\), by Proposition~1(a) and the left
side of (H1). So only the zero-cost challenger prepares, and entry stays at the floor \(\rho\) after every order, at
both strengths.

The investor still never trades at \(r_0\), and it still trades fully at \(r_1\), where the trading bound with
entry fixed at \(\rho\) exceeds the trading cost (Appendix~A.6 of the paper). The price at \(r_1\) carries
information, but nobody who decides on entry reads it.

\subsection{What the two controls show together}

Each control removes one link of the chain in Proposition~2.
\begin{itemize}
\item Panel~B removes the change in information and keeps the change in strength. Entry does not rise. So the
  stronger incumbent, by itself, does not create the rise.
\item Panel~C keeps the trade that the stronger incumbent creates, and it cuts the link from the price to the
  challenger. Entry does not rise. So trade, by itself, does not create the rise either.
\end{itemize}
Only Panel~A has both links: the stronger incumbent makes informed trade pay, the trade puts information into the
price, and the challenger reads it. The paper concludes that the two controls locate the reversal in the information
that trading creates (Section~4.2 of the paper). Incumbent strength changes that information, and the entry response
to it exceeds the direct deterrence of a stronger incumbent at fixed information.

A picture to keep in mind, which is the guide's own use of the table, splits the rise from \(r_0\) to \(r_1\) into two
moves. First, at the weak incumbent, switch from the silent price of Panel~A to the full-order information of
Panel~B: entry rises from \(\rho\) to \(\rho+(1-\rho)(b+\lambda/2)\), the response to information. Second, at
full-order information, strengthen the incumbent from \(r_0\) to \(r_1\): entry does not change, so deterrence at
fixed information is zero over this step. The two moves together give the rise of Panel~A. Deterrence at fixed information becomes positive only above the
entry ceiling, where it removes the recruitment after a purchase. That is the fall from \(r_1\) to \(r_2\).

\begin{checkyourself}
\label{q:example-hidden}
In Panel~C the challenger decides on its prior.
(a)~Show that the investor at \(r_1\) still buys in class \(H\) and sells in class \(L\), although no purchase recruits.
(b)~Write the three prices at \(r_1\), and show that they differ.
(c)~Explain why entry is still \(\rho\).
\end{checkyourself}

% ====================================================================
\begin{chaptersummary}
\begin{itemize}
\item The example has one value in each class, a reserve below both values, and an incumbent uniform on \([0,r]\)
  with \(\ell<r<h\); a larger \(r\) is a stronger incumbent, and class \(H\) always wins.
\item Its auction objects have closed forms: \(g_H=h-r/2-p^2/(2r)\), \(g_L=(\ell^2-p^2)/(2r)\) and
  \(\DT=(r-\ell)^2/(2r)\); both profits fall strictly and the spread rises strictly with \(r\), and
  \(\DT+g_H-g_L=h-\ell\).
\item The threshold \(\tau(r)=[2rc_H-(\ell^2-p^2)]/(2rh-r^2-\ell^2)\) rises slowly, by less than \(1/(2\gamma)\) per unit
  of strength, because class \(H\)'s payment rises by less than one half per unit.
\item At the prior \(\tfrac12\), (H1)--(H3) hold exactly when the top posterior lies between \(\tau(r_1)\) and
  \(\tau(r_2)\); they depend on the slot only through \(u\), and they hold on an open set of parameters.
\item The entry ceiling \(r_C\), where \(\tau(r_C)=M(\tfrac12)\), lies just below \(r_2\); with a high cost above
  \((h^2-p^2)/(2h)\), class \(H\) outbids every uniform incumbent against which a price recruits.
\item At fixed information, a stronger incumbent can only lower entry, by the probability that the cost lies between
  the two gross profits; in the slot this needs a positive spread in each economy.
\item With full orders held fixed, entry is flat from \(r_0\) to \(r_1\); with the price hidden, entry stays at
  \(\rho\); so the rise comes from the information that trading puts into the price.
\end{itemize}
\end{chaptersummary}

% ====================================================================
\answersection

\answer{\ref{q:example-reserve}}{%
(a)~In Table~\ref{tab:example-cases}, the first two rows give the same price for both classes: \(p\) when \(R<p\),
and \(R\) when \(p\le R<\ell\). Only the third row differs: class \(H\) pays \(R\), and with class \(L\) the incumbent
pays \(\ell\). In that row \(R>\ell>p\), so the reserve does not bind for either class. The reserve enters both
classes' proceeds in the same way and cancels in the difference:
\begin{align*}
t_H-t_L&=\frac1r\int_\ell^r(y-\ell)\,dy \why{only the third row differs}\\
&=\frac{(r-\ell)^2}{2r} \why{the integral}
\end{align*}
(b)~\(g_L=(\ell^2-p^2)/(2r)\) falls as \(p\) rises, and it tends to zero as \(p\) approaches \(\ell\). In words, class
\(L\) earns \(\ell-\max\{p,R\}\) only when \(R<\ell\). This margin is at most \(\ell-p\), so a reserve near \(\ell\)
leaves class \(L\) almost nothing.}

\answer{\ref{q:example-ceiling}}{%
(a)~In the band, \(\tau(r_1)<M(\tfrac12)<\tau(r_2)\). The threshold is continuous and rises strictly on the path
(Section~\ref{sec:example-shape}). The intermediate value theorem gives a strength \(r_C\) in \((r_1,r_2)\) with
\(\tau(r_C)=M(\tfrac12)\), and strict monotonicity makes it unique.

(b)~\(M(\tfrac12)=(1+u)/(2+u)=1-1/(2+u)\) rises strictly with \(u\). The threshold curve does not depend on the
slot. A higher level \(M\) meets an increasing curve further to the right, so \(r_C\) moves to a higher strength. In words,
more informed presence relative to liquidity makes a purchase stronger news, and stronger news
recruits against stronger incumbents.

(c)~For \(u_{\rm hi}\le u<u_k\), \(M\ge\tau(r_2)\), so \(r_C\ge r_2\). Orders are still full at \(r_2\), because the
trading bound forces them (Corollary~1(b)). A purchase at \(r_2\) lifts the belief to \(M\ge\tau(r_2)\) and recruits
under the default tie rule. Entry at \(r_2\) is then \(\rho+(1-\rho)(b+\lambda/2)\) instead of \(\rho\), and the fall
of entry from \(r_1\) to \(r_2\) disappears.}

\answer{\ref{q:example-width}}{%
The fundamental theorem of calculus writes the rise of \(\tau\) as the integral of its slope.
\begin{align*}
\tau(r_2)-\tau(r_1)&=\int_{r_1}^{r_2}\tau'(r)\,dr \why{fundamental theorem of calculus}\\
&<\int_{r_1}^{r_2}\frac{1}{2\gamma(r)}\,dr \why{\(\tau'<1/(2\gamma)\)}\\
&\le\int_{r_1}^{r_2}\frac{1}{2\gamma(r_2)}\,dr \why{\(\gamma(r)\ge\gamma(r_2)>0\) for \(r\le r_2\)}\\
&=\frac{r_2-r_1}{2\gamma(r_2)} \why{a constant integrand}
\end{align*}
The third line uses \(\gamma'=-\DT'<0\): the profit gap falls with \(r\), so it is smallest at \(r_2\), and its
reciprocal is largest there. The bound \(\tau'<1/(2\gamma)\) needs \(0\le\tau\le1\) on \([r_1,r_2]\), which holds in the
example.

In words: per unit of strength, gross profit at the threshold falls by less than one half, and the belief must
rise by that fall divided by the slope \(\gamma\) of gross profit in the belief. A steep line, that is, a large
profit gap, needs only a small move of the belief, so the two ends of the band lie close together.}

\answer{\ref{q:example-control}}{%
With full orders the posteriors are \((m,\tfrac12,M)\) at every strength. A sale and no order never recruit, because
\(\tau>\tfrac12>m\) on the path. A purchase recruits at a strength \(s\) exactly when \(s\le r_C\). By
Proposition~A.2 with the two-point cost, the probability is \((1-\rho)\) times the probability of an order that
recruits at \(r\) and not at \(r'\).
\begin{itemize}
\item[(a)] \(r<r'\le r_C\): a purchase recruits at both strengths. No order switches, so the probability is zero.
  Gross profit at \(M\) falls, but it still covers \(c_H\) at \(r'\). Entry stays at \(\rho+(1-\rho)(b+\lambda/2)\).
\item[(b)] \(r\le r_C<r'\): a purchase recruits at \(r\) and not at \(r'\). The probability is
  \((1-\rho)\Pr(X=1)=(1-\rho)(b+\lambda/2)\). Entry falls from \(\rho+(1-\rho)(b+\lambda/2)\) to \(\rho\).
\item[(c)] \(r_C<r<r'\): no order recruits at either strength. The probability is zero, and entry stays at \(\rho\).
\end{itemize}
So entry with information held fixed is a step: flat, then one drop at the ceiling, then flat. It never rises.}

\answer{\ref{q:example-hidden}}{%
(a)~The challenger ignores the price, so entry is \(e(x)=\rho\) after every order in every candidate. The trading
returns of the two correct orders then satisfy
\begin{align*}
K_H&=\rho\,\Delta_1\,[1-\mu(1)] \why{\eqref{eq:trading-returns} with \(e(1)=\rho\)}\\
&\ge\rho\,\Delta_1\,[1-M] \why{\(\mu(1)\le M\), Lemma~A.1(b)}\\
&\ge\rho\,\kappa\,\Delta_1>k \why{definition of \(\kappa\); right side of (H2)}
\end{align*}
and in the same way \(K_L=\rho\Delta_1\mu(-1)\ge\rho\, m\,\Delta_1\ge\rho\kappa\Delta_1>k\). So in every candidate a
purchase beats abstention in class \(H\), and a sale beats abstention in class \(L\) (Lemma~A.3(b)). Wrong-signed
orders lose (Lemma~A.3(a)). So the investor trades fully, and no recruitment is needed for this.

(b)~With full orders, the posteriors are \((m,\pi,M)\). Write \(w_L=t_{L,1}-t_{0,1}\). The prices are
\(P(x)=t_{0,1}+\rho\,[w_L+\Delta_1\mu(x)]\), so
\begin{align*}
P(0)-P(-1)&=\rho\,\Delta_1\,(\pi-m) \why{\(t_{0,1}\) and \(\rho w_L\) cancel}\\
P(1)-P(0)&=\rho\,\Delta_1\,(M-\pi) \why{same steps}
\end{align*}
Both differences are positive, because \(\rho>0\), \(\Delta_1>0\) and \(m<\pi<M\). So the three prices differ, and the
price moves with the order.

(c)~The challenger does not see the price, so its belief stays at the prior \(\pi\). Gross profit there is
\(B_{r_1}(\pi)\le B_{r_0}(\pi)<c_H\), by Proposition~1(a) and the left side of (H1). A high-cost challenger never
prepares, and only the zero-cost challenger does. Entry is \(\rho\) after every order. The information is in the
price, but nobody who decides on entry reads it.}
